\section{Introduction}\label{sec:intro}

Video summarization reduces long audiovisual collections to selected keyframes,
shots or segments~\citep{truong2007video_abstraction,apostolidis2021survey}.
When several sources cover the same topic, selection must also handle
cross-source redundancy~\citep{li2010videommr,panda2017collaborative,
wu2020mvsgcn}. A multi-video pipeline therefore combines several criteria: a
segment can be relevant yet redundant, visually informative yet technically
poor, or temporally misplaced.

HSMVideoSumm makes those criteria inspectable. A study with human summarizers of
multi-source news identified five recurrent strategies---removing introductory
passages, filtering subjective content, detecting redundancy, selecting for
visual quality and ordering chronologically---and implemented them as separate
stages~\citep{barbieri2021content,barbieri2021tese}. The present study concerns
only the fourth stage. Segmentation, subjectivity filtering, redundancy
detection and ordering remain unchanged. Once the pipeline treats segments as a
redundant cluster, visual quality acts as a tie-breaker; this conditional design
does not imply that cluster members have identical semantic utility.

The scientific issue is the distinction between a criterion and the operation
that implements it. The criterion is human-grounded because it was observed in
the summarization study. Its implementation is a later engineering choice:
Scale-Invariant Feature Transform (SIFT) descriptors are quantized through
Bag-of-Visual-Words (BoVW), weighted with term frequency--inverse document
frequency (TF-IDF), and summed~\citep{lowe2004sift,sivic2003videogoogle,
csurka2004bovw,salton1988tfidf}. SIFT, BoVW and TF-IDF were developed for
matching, retrieval and categorization, not for predicting perceptual quality.
Calling their aggregate a quality score does not establish construct validity.

No-reference image quality assessment (NR-IQA) provides alternatives because
segments from different broadcasters have no common undistorted reference. The
available families include natural-scene-statistics, supervised convolutional
and transformer regressors, and vision--language models
~\citep{mittal2012brisque,mittal2013niqe,ke2021musiq,zhang2023liqe,
agnolucci2024qualiclip}. The question is not whether a newer model produces a
higher number, but whether the deployed operation responds to the construct it
names, whether that behaviour can be explained, and whether replacing the
operation changes the pipeline output.

The study makes three contributions. First, it derives an exact factorization
that isolates the descriptor-count and mean-IDF-weight channels of the deployed
score, then measures which channel dominates across image databases, broadcast
segments and 19 convolutional neural network (CNN) extractor substitutions.
Second, it combines human-score
association with within-reference degradation sensitivity and a constructed
experiment containing 800 paired decisions between an unmodified signal and a
deterministically degraded counterpart. Third, it applies a common screening
rule and substitutes the complete visual-quality operation in 42 redundancy
clusters while holding every other pipeline stage fixed, measuring both changed
selections and model-reference scores of the selected segments. Training-database overlap,
detectable-keypoint conditioning, alternative pooling, overlay strata,
predictions recorded in a commit whose metadata predates the experiment and
computational details are auxiliary diagnostics reported in the supplementary
material.

The resulting claim is deliberately bounded: the factorization explains the
observed count dominance and failure of the evaluated BoVW+TF-IDF
configurations. It is not a theorem that every possible visual vocabulary,
normalization or extractor must fail. Likewise, the application study separates
Q-Align, which defines the in-domain screening floor, from DOVER, which is used
only for output sensitivity; neither model substitutes for human preference
judgments.
